\section{Gate Details, Test Bed and Ensemble Variant}
\label{sup:testbed}
\subsection{Inverse Problem: Pose Graph Optimisation with Depth Maps}
\label{sec:method:inverse}
The poses and the alignments are estimated jointly,
\begin{equation}
\begin{split}
  \min_{\{\mathbf{T}_i\},\{\mathbf{a}_{i}\}}\;
  \sum_{(i,j)\in\mathcal{E}}\Bigl[&\rho_{g}\bigl(\|\mathbf{r}^{\mathrm{geo}}_{ij}\|^{2}_{\Sigma_g}\bigr)\\
  &+g_{i}\,g_{j}\,
  \rho_{d}\bigl(\|\mathbf{r}^{\mathrm{dep}}_{ij}\|^{2}_{\Sigma_d}\bigr)\Bigr],
\end{split}
  \label{eq:pgo}
\end{equation}
where $\mathbf{r}^{\mathrm{geo}}_{ij}$ are epipolar or reprojection residuals of the
feature matches between views $i$ and $j$, $\rho_g$ and $\rho_d$ are robust kernels
(a quadratic kernel when no robust baseline is used), and the depth-induced residual
back-projects a pixel $u$ of view $i$ with its aligned depth, transforms it into view
$j$, and compares its depth there with the depth of view $j$:
\begin{equation}
\begin{split}
  r^{\mathrm{dep}}_{ij}(u)&=\tilde D_{j}\bigl(\pi_j(\mathbf{q})\bigr)-q_z,\\
  \mathbf{q}&=\mathbf{T}_j^{-1}\mathbf{T}_i\,\pi_i^{-1}\bigl(u,\tilde D_{i}(u)\bigr).
\end{split}
  \label{eq:rdep}
\end{equation}
The weights $g_ig_j$ are where the gate acts: a depth constraint is built only from two
kept depth maps. These constraints are pairwise because the pose graph is, but they are
built after the gate, only from kept depth maps, and they are not part of the gate's cost
(Section~\ref{sec:method:cost}). The back-end is a differentiable least-squares solver \cite{opt_theseus} and the surface is
renderable, so the depth residual of a view against the rendered surface,
$\tilde D_i-\hat S_i(\mathbf{T}_i)$, can serve as an additional, linear-cost constraint on
$\mathbf{T}_i$; in this paper the gate is training-free and the back-end is used as an ordinary
solver \cite{opt_g2o,opt_factorgraphs}. The differentiable, renderable optimiser is
the test bed the gate plugs into, not the framing. Its closest relative of the gate is
\emph{soft down-weighting}: replacing $g_i\in\{0,1\}$ by a weight $w_i\in[0,1]$ estimated
inside the optimisation from the view's residual. This is a baseline (H5): it reduces the
influence of a view but neither excludes it nor states a false-alarm or missed-detection
probability. 


\subsection{Test Bed: Differentiable, Renderable Pose Graph}
\label{sec:method:pipeline}
The test bed runs five steps. (i)~\emph{Depth and matches:} each view is processed by the
depth model, and feature matches are computed for view pairs. (ii)~\emph{Provisional poses:}
a pose graph with the geometric constraints only is solved with a robust kernel, to give
poses that are good enough to render a surface at each view. (iii)~\emph{Gate:}
Section~\ref{sec:method:gate}: fuse the shared surface in information form, subtract each view, and test every depth map against
its rendering. (iv)~\emph{Pose graph optimisation:} \eqref{eq:pgo}, with gated depth
constraints and, optionally, a robust kernel, a soft per-view weight (baseline) or the
render-based constraint of Section~\ref{sec:method:inverse}. (v)~\emph{Fused output:} the primitives are
re-fused from the kept depth maps at the refined poses; its rendered depth is the fused depth,
and a truncated signed distance function \cite{fus_curless,fus_kinectfusion} of the same
kept depth is a secondary check. The gate may be repeated once with the refined poses.
Because the gate relies on provisional poses, errors in those poses can corrupt both the
surface and the gate; the repetition, and the comparison of gating before versus after the
pose graph (H5), test this. The baseline families replace or complement step (iii)
(Section~\ref{sec:setup:baselines}).

\subsection{Ensemble Variant (Ablation)}
\label{sec:method:ensemble}\label{sup:ensemble}
The ensemble ablation gives each view $M$ depth maps $D_{m,i}$ from $M$ models and asks two
questions: whether models that share training data fail on the same views, and whether
voting across models breaks where the gate does not. Cross-model voting takes, at each
pixel, the median (or majority) of the $M$ aligned depth maps and flags a view when the
models disagree beyond their noise; it needs no other view. The ensemble gate tests each of
the $M$ depth maps of a view against the same surface rendered at the view, so that
$g_{m,i}$ is defined per model and a view is dropped when every model is rejected or when
the kept models disagree. Both are evaluated by H10 (Section~\ref{sec:setup:hyp}): they
should differ most on views where several models fail together, since cross-model voting then
agrees with itself, whereas the surface is built from other views.


\subsection{Cost of the Gate}
\label{sup:cost}
A pairwise gate warps the depth maps of the neighbouring views into the view under test
(Appendix~\ref{app:pairwise}). With $P$ pixels per view, mean co-visibility degree $\bar d$
and a warp cost $c_w$ per pixel, its cost is $C_{\mathrm{pair}}\approx N\bar d\,P\,c_w$,
which grows with the number of view pairs when overlap is dense ($\bar d\propto N$). The
shared-surface gate builds the primitives of each view once, at a cost $c_b$ per pixel, and
renders every view once at a cost $c_r$ per pixel; the leave-view-out is an $O(1)$
subtraction per primitive. Its cost is $C_{\mathrm{surf}}\approx N\,P\,(c_b+c_r)$, linear in $N$,
with no optimisation iterations and no refit per fold. Association adds a lookup per pixel. The
surface therefore wins when overlap is dense: the break-even degree is
$\bar d^{*}\approx(c_b+c_r)/c_w$, and for sparse view graphs the pairwise gate can be cheaper.
The primitive count grows with the scene rather than being fixed, which is the price of dropping
a fixed budget. Depth inference adds $N$ model evaluations per scene, which
can be cached and which do not depend on the gate; with $M$ models the ensemble ablation
multiplies this term by $M$. H7 measures the scaling and the break-even point.


\subsection{Minimum Detectable Bias Table}
\begin{table}[t]
\caption{\textbf{Minimum detectable bias} (in $\sigma$) for a bias on a fraction $f$ of $n=1000$
independent pixels, $\alpha=0.05$, $\beta_{\mathrm{md}}=0.1$. Matched: region known; energy: unknown. $f=1$ is absorbed by the offset.}
\label{tab:mdb}
\centering
\small
\setlength{\tabcolsep}{4pt}
\begin{tabular}{lccc}
\toprule
$f$ & Matched & Scan, $R{=}100$ & Energy \\
\midrule
0.05 & 0.47 & 0.69 & 1.71 \\
0.2  & 0.26 & 0.38 & 0.93 \\
0.5  & 0.21 & 0.30 & 0.75 \\
1    & $\infty$ & $\infty$ & $\infty$ \\
\bottomrule
\end{tabular}
\end{table}

\subsection{Relation to Leave-One-Out and Solution-Separation Tests}
\begin{proposition}[Relation to leave-one-out and solution-separation tests]
\label{prop:loo}
Let source $k$ measure a state $x$ directly with noise covariance $R_k$, let
$\hat x_{-k}$ be the estimate without it with covariance $P_{-k}$, and let
$\nu=z_k-\hat x_{-k}$ with $S=R_k+P_{-k}$. Then the solution-separation statistic
$d^{\top}\mathrm{Cov}(d)^{-1}d$, with $d=\hat x_{\mathrm{all}}-\hat x_{-k}$,
equals the leave-one-out innovation statistic $\nu^{\top}S^{-1}\nu$. For Gaussian
PDFs the negative log-evidence of $k$ under the consensus,
$-\log\!\int p_C f_k=\tfrac12\nu^{\top}S^{-1}\nu+\tfrac12\log\det(2\pi S)$, is the
leave-one-out statistic plus a covariance-volume term.
\end{proposition}
(Proof in Appendix~\ref{app:proof}.)
\begin{remark}
Under equal covariances and Gaussian residuals, ranking depth maps by this evidence is
therefore the same as ranking by the leave-one-out chi-square statistic, and the
evidence criterion can add value only where noise is heteroscedastic (as depth noise is),
residuals are non-Gaussian or multimodal, or occlusion and grazing geometry break the
chi-square assumptions. We test exactly this in H3 (Section~\ref{sec:setup:hyp}) and treat
a tie as a reportable outcome.
\end{remark}


\todofig{pipeline}{Test bed: RGB images; one metric depth model per image and feature matches; provisional poses; shared Gaussian surface fused in information form, with leave-view-out by subtraction; \emph{view-level gate} against the rendered surface (contribution); pose graph optimisation with gated depth factors; surface of the kept depth maps as the fused output (posterior, protection level, view count).}

